\section{The language \Lang{}}
\label{sec:language}

\Lang{} is the source side of the theorem. It is meant to be recognizable:
a reader who knows POSIX-style threads should find no surprise in its
constructs. Its semantics, however, is fixed in detail, because an exact
equivalence leaves no room for ``implementation-defined''.

\subsection{Syntax}

Figure~\ref{fig:syntax} gives the syntax (\lean{L.Stmt}, \lean{L.Expr},
\file{L/Syntax.lean}). A program has a list of procedures, a \kw{main} body
with one parameter that receives the program input, and a number of
\emph{statically known} mutexes, which exist and are unlocked when the
program begins. Expressions are pure and read thread-local variables only.
Shared variables are read by \kw{load} and written by \kw{store}, one
access per statement.

\begin{figure}[t]
\small
\[
\begin{array}{@{}r@{\;}c@{\;}l@{}}
\mathit{lit} &::=& \kw{null} \mid \mathit{bool} \mid \mathit{int} \mid \mathit{string} \\[2pt]
e &::=& \mathit{lit} \mid x \mid \kw{staticMutex}(i) \mid \kw{cvOf}(e, i)
      \mid e + e \mid e = e \mid e < e \mid \lnot e \\[2pt]
s &::=& \kw{skip} \mid s ; s \mid x := e \mid x := \kw{load}\ g \mid \kw{store}\ g := e
      \mid \kw{if}\ e\ \kw{then}\ s\ \kw{else}\ s \mid \kw{while}\ e\ \kw{do}\ s \\
  &\mid& x := p(\vec{e}) \mid \kw{return}\ e
      \mid x := \kw{fork}\ p(\vec{e}) \mid \kw{join}(e) \mid \kw{cancel}(e) \\
  &\mid& x := \kw{newMutex} \mid \kw{lock}(e) \mid \kw{unlock}(e)
      \mid \kw{wait}(e_{\mathit{cv}}, e_m) \mid \kw{signal}(e) \mid \kw{broadcast}(e) \\
  &\mid& \kw{sleep}(e) \mid x := \kw{recv}(c)
      \mid \kw{select}\ \{\, x := \kw{recv}(c) \to s \;\mid\; \kw{after}(e) \to s \,\} \\
  &\mid& \kw{emit}(e) \mid \kw{throw}\ e \mid \kw{try}\ s\ \kw{catch}\ (x)\ s
      \mid \kw{exit}(e) \mid \kw{fail}(e)
\end{array}
\]
\caption{The syntax of \Lang{}. $x$ is a thread-local variable, $g$ a shared
variable, $p$ a procedure, $c$ an external channel, and $i$ an index.}
\label{fig:syntax}
\end{figure}

Values are the literals plus three kinds of handles: thread handles, mutex
handles, and condition variable handles. A condition variable is named by
its mutex and an index, so it belongs to exactly one mutex.

\Lang{}'s observable alphabet \ObsL{} has six events (\lean{L.Event}):
$\ev{begin}(v)$, the program starts with input $v$; $\ev{in}(c, v)$, a
message arrives on external channel $c$; $\ev{refused}(c, v)$, a message
arrives after the program ended; $\ev{out}(a)$, an emitted effect is
delivered; $\ev{exit}(v)$; and $\ev{fail}(s)$.

\subsection{Semantics}
\label{sec:l-semantics}

The semantics is a small-step transition system (\lean{L.step},
\file{L/Semantics.lean}). A configuration holds a status (not started,
running, closed), a clock, the live threads, the shared variables, the owner
of each mutex, the FIFO queue of each condition variable, the queued
messages of each external channel, the cancel flags, the threads that ended
and were not joined, and, after the end, the emits that were issued and not
yet delivered. A label is an idle clock tick, an environment request, the
next action of one thread, the internal action \kw{fire}, or the delivery or
drop of a late emit.

The definition has two layers. \file{L/ThreadLocal.lean} says what a
statement does to its own thread and which atomic action it asks for, as a
value of type \lean{L.Intent}. \file{L/Semantics.lean} performs that action
on the shared state. The translation reuses the first layer verbatim
(\S\ref{sec:translation}), which is what keeps the two sides aligned.

The semantic choices below are the ones that matter for the equivalence.
Each matches what the Dex core does.

\begin{description}[leftmargin=1.2em,itemsep=3pt]
\item[Atomic actions.] Each single read or write of a shared variable, each
  issue of an \kw{emit}, and each synchronization operation is one atomic
  action. Local computation is invisible. There is no atomic
  read-modify-write: $g := g + 1$ is a \kw{load} and a \kw{store}, and two
  threads can interleave them.
\item[Mutexes.] Threads blocked in \kw{lock}$(m)$ acquire $m$ in FIFO order
  of blocking. A mutex that does not exist is never free.
\item[Condition variables.] \kw{wait}$(\mathit{cv}, m)$ atomically releases
  $m$ and appends the thread to $\mathit{cv}$'s FIFO queue. \kw{signal}
  wakes the head of the queue, \kw{broadcast} wakes all. A woken thread
  re-locks $m$, in FIFO order with other lockers. There are no spurious
  wakeups.
\item[Cancellation] is cooperative and deferred. \kw{cancel}$(t)$ sets
  $t$'s flag. Thread $t$ observes the flag at a cancellation point
  (\kw{sleep}, \kw{recv}, \kw{select}, \kw{join}), and only if that point
  would block; it then ends without releasing anything it holds, as in POSIX
  without cleanup handlers. \kw{lock} and \kw{wait} are not cancellation
  points. A canceled thread is joinable like any other ended thread.
\item[Emit is asynchronous.] \kw{emit}$(a)$ is \emph{issued} as one atomic
  action. The event $\ev{out}(a)$ happens later, after any interleaving of
  other threads, and the thread continues only after it. If the program
  ends between the issue and $\ev{out}(a)$, the effect may still be
  delivered after $\ev{exit}$, or never.
\item[Exit.] \kw{exit}$(v)$, \kw{fail}$(s)$, or the return of \kw{main}
  ends the program at once: every other thread stops. A thread that ends
  otherwise just ends, and a program whose threads all end without
  \kw{exit} is stuck forever. Input that arrives after the end is
  $\ev{refused}$.
\item[Join.] A thread that ends leaves one join token. \kw{join}$(t)$ takes
  the token if it is there; otherwise it observes its own cancel flag; and
  otherwise it blocks. A second \kw{join} of the same thread blocks
  forever.
\item[Exceptions] are data. \kw{throw} transfers control to the innermost
  handler on the thread's stack. An uncaught exception whose value is a
  string $s$ fails the program with $\ev{fail}(s)$.
\item[Halting.] \Lang{} has no static types. A statement whose operand has
  the wrong type, and an uncaught exception that is not a string, halt the
  thread: it takes no further action, never ends, and releases nothing.
\end{description}

\subsection{The timing model}
\label{sec:timing}

\Lang{} has an internal clock in microseconds, advanced by idle ticks. A
timed operation's deadline is the clock value when it is armed plus a whole
number of seconds. Dex computes fire times the same way, from commit
timestamps, and those strictly increase from commit to commit. So two runs,
one on each side, can carry the same deadlines only if the same number of
clock-stamped steps lies between any two arming points. If the counts
differ, the order of two deadlines can differ, and a \kw{select} makes that
order observable.

\Lang{} therefore has a \emph{fixed timing model}: each kind of statement
takes a fixed number of \emph{stamped actions}, and the clock must strictly
increase between any two stamped actions. The counts
(Table~\ref{tab:timing}) are the numbers of Dex commits of the translation.

\begin{table}[t]
\centering
\small
\begin{tabularx}{\textwidth}{@{}LP{5.6cm}@{}}
\toprule
\Lang{} step & Stamped actions \\
\midrule
$\ev{begin}(v)$ & 1, then 2 more before \kw{main}'s first statement \\
$s ; s$ and \kw{skip} & 0 \\
A statement that does not block (assignment, \kw{load}, \kw{store},
  \kw{if}, each test of \kw{while}, call, \kw{return}, \kw{fork},
  \kw{newMutex}, \kw{unlock}, \kw{signal}, \kw{broadcast}, \kw{emit},
  \kw{throw}, entering \kw{try}, \kw{cancel}, \kw{exit}, \kw{fail}); the end
  of a \kw{try} body or a procedure body & 2: begin, complete \\
\kw{lock}, \kw{recv}, \kw{join} & 4: begin the wait, register it, begin the
  continuation, complete it \\
\kw{wait}$(\mathit{cv}, m)$ & 10: 2 to release $m$ and enqueue, 4 to wait
  for the signal, 4 to re-lock $m$ \\
\kw{sleep}, \kw{select} with \kw{after} & 6: 2 for Enter, then 4, of which
  the register is Arm \\
$\ev{in}(c, v)$ while the program runs & 1 \\
A batch of $k \ge 1$ consecutive \kw{fire} actions & 1 \\
$\ev{out}(a)$; delivery or drop of a late emit; $\ev{refused}(c, v)$ & 0 \\
\bottomrule
\end{tabularx}
\caption{The timing model of \Lang{}. A shared read happens at begin, an
\kw{emit} is issued at begin, and every other atomic action happens at
complete or at register.}
\label{tab:timing}
\end{table}

This was an explicit decision of the model's owner (row L1 of the
specification's review log), with this rationale: real programs cannot rely
on timing races. A program whose outcomes do not depend on the order of two
deadlines means the same under this timing model as under any other. The
theorem adopts the one model that is aligned with Dex, which makes the
equivalence exact without narrowing what it covers in practice. It is like
counting one time unit per instruction in algorithm analysis. The theorem
does \emph{not} say that a program is equivalent to itself under another
cost model.

In Lean the model is the phase of a thread (\lean{L.Phase}): a statement
that does not block goes \lean{ready} $\to$ \lean{running} $\to$
\lean{ready}; an \kw{emit} goes \lean{ready} $\to$ \lean{emitting} $\to$
\lean{emitted} $\to$ \lean{ready}; a blocking statement goes \lean{ready}
$\to$ \lean{waitStarted} $\to$ \lean{blocked} or \lean{woken} $\to$
\lean{resuming} $\to$ \lean{ready}.

\subsection{Timers}
\label{sec:timers}

A timed operation takes two actions. \emph{Enter} gives it the next timer
number $n$ from one program-wide counter. \emph{Arm}, later, sets its
deadline to the current clock plus the duration and blocks the thread. The
timer completes only by the internal action \kw{fire}, which is enabled at
any time no earlier than the deadline and possibly much later. \kw{fire}
fires only the \emph{least} due timer in the order (deadline, $n$). A batch
of consecutive \kw{fire} actions shares one clock stamp, as one
\rulename{TIMER-FIRE} commit fires a prefix of the due timers. A blocked
timed operation is decided by whichever comes first: its \kw{fire}, the
arrival of its message, or a cancel. A timer is never complete when it is
armed, even with duration zero.

\subsection{Scope decisions}
\label{sec:scope}

Three more decisions bound the theorem. Each is recorded in the review log
of the specification with its date and its reason.

\paragraph{Well-formed programs (L2).}
\kw{wait}, \kw{signal}, and \kw{broadcast} require the condition variable's
mutex to be held, and \kw{unlock} requires the mutex to be held by the
caller. \Lang{} gives no behavior to a thread that is about to break a
requirement: the thread has no step. The theorem covers the programs in
which this never happens.

\begin{Verbatim}
def Thread.BreaksMutexRequirement (config : Config) (thread : Thread) : Prop :=
  thread.phase = .ready ∧ thread.timer = none ∧
  ∃ mutex, requiredMutex thread.core thread.code = some mutex ∧
    config.mutexOwners.get? mutex ≠ some (some thread.core.tid)

def BreaksMutexRequirement (config : Config) : Prop :=
  ∃ thread, thread ∈ config.threads ∧ thread.BreaksMutexRequirement config

def Program.IsWellFormed (program : Program) : Prop :=
  ∀ config, (system program).Reachable config → ¬ BreaksMutexRequirement config
\end{Verbatim}

\lean{Program.IsWellFormed} is a property of the source program under
\Lang{}'s semantics alone, quantified over every input and every schedule,
like data-race freedom. The translation plays no part in it.

\paragraph{Unbounded counters (L3).}
Step numbers and timer numbers are natural numbers in the model. The real
server stores a Step number in 32 bits. The order of timers with equal fire
times is derived from the byte order of encoded timer IDs, and that
derivation is exact for fewer than $2^{32}$ Step executions of one type per
run. The theorem says nothing about a run that exceeds the limit.

\paragraph{Unforgeable handles and literal-only inputs (L5, L6).}
The proof itself found these two points. In an earlier version a literal
could be any value, including the handle of a mutex that \kw{newMutex} had
not created yet. Consider
\begin{quote}\small
$\mathit{locker}()$: \kw{lock}$(h)$; \kw{emit}$(1)$ \qquad
\kw{main}: \kw{fork} $\mathit{locker}()$; \kw{sleep}$(0)$;
\kw{newMutex}; \kw{recv}$(9)$
\end{quote}
where $h$ is the forged handle of the mutex that \kw{main} creates. In
\Lang{} the locker blocks forever, because creating a mutex does not serve
blocked lockers. In Dex the first token that the translated \kw{newMutex}
publishes wakes the waiting Step, and every fair run emits $\ev{out}(1)$.
The liveness statement was false for this well-formed program. The fix is
at the source: literals are only null, Booleans, integers, and strings
(\lean{L.Literal}); a mutex handle comes only from \kw{newMutex} or
\kw{staticMutex}; a thread handle comes only from \kw{fork}. The lemmas
\lean{evalExpr_mutex_origin} and \lean{evalExpr_thread_origin} state this.

The counterexample then came back through the environment: the program
input was an arbitrary value, so it could be the forged handle. The second
decision closes that door: the environment sends plain data only. A request
carries a \lean{L.Literal}, and \lean{literal_isNotHandle} proves that no
literal is a handle. After both decisions the counterexample cannot be
written down, which is a stronger guarantee than a side condition that
excludes it.
